\documentclass[a4paper,12pt]{article}
\usepackage[finnish]{babel}
\usepackage[utf8]{inputenc}
\usepackage[T1]{fontenc}
\usepackage{lmodern}
\usepackage{eurosym}
\usepackage[pdftex,colorlinks]{hyperref}
\renewcommand{\thesubsection}{\arabic{subsection}}
\title{Etelä-Espoon sosialidemokraattinen työväenyhdistys r.y.}
\author{EED}
\date{22.10.2026}
\begin{document}
\maketitle
\tableofcontents
\section*{Johtokunnan kokous}
\subsection*{Paikka ja aika}
Ravintola Tillinkulmassa, torstaina 22.10.2026 kello 17:30.
\section*{Läsnä}
Markku Sistonen, puheenjohtaja \\
%Arja Kivimäki, rahastonhoitaja \\
Martti Hellström, tiedotusvastaava \\
Pauliina Sistonen, tiedotusvastaava \\
%Jonne Hänninen \\
Matti Ikonen \\
Mikko Nummelin, sihteeri
\section*{Pöytäkirja}
\subsection{Kokouksen avaus}
Puheenjohtaja avasi kokouksen kello 17:40.
\subsection{Kokouksen laillisuus ja päätösvaltaisuus}
Kokous todettiin lailliseksi ja päätösvaltaiseksi.
\subsection{Kokouksen ääntenlaskijoiden valinta}
Kokouksen ääntenlaskijoiksi valittiin Martti Hellström ja Pauliina Sistonen..
\subsection{Jäsenasiat}
Uudeksi jäseneksi valittiin Anneli Tolvanen.
\subsection{Syyskokouksen valmistelu}
Arja Kivimäki lähetti talousarvion ja Markku Sistonen tekee toimintasuunnitelman. Jo hyväksyttyjä sääntöjä joudutaan muuttamaan suhteellisen vaalitavan vaatimisen osalta. Lisäksi kokouksen on hyväksyttävä vuoden 2025 johtokunnan vastuuvapaus. Päätetään, pyydetäänkö uusi johtokunta valmistelemaan Falklammen käytön tulevaisuutta, sisältäen mahdollisuuden Falklammen myyntiin tai käytön laajentamiselle.
\subsection{Syksyn tapahtumat}
\begin{itemize}
\item{Johtokunnan kokous} to 22.10. kello 17:30, ravintola Tillinkulmassa
\item{Syyskokous} viikolla 47
\item{Puurojuhla} ke 2.12 kello 18, Ravintola Tillinkulma
\end{itemize}
\subsection{Muut asiat}
\subsection{Kokouksen päättäminen}
Puheenjohtaja päätti kokouksen kello 18:30.
\end{document}
